% sec:gr-stress-energy — Stress-Energy and Matter Coupling
\section{Stress-Energy and Matter Coupling}
\label{sec:gr-stress-energy}

The left-hand side of the Einstein equations encodes geometry; the
right-hand side encodes the matter and energy that source it.  But what,
concretely, does the stress-energy tensor $\stressenergy$ look like for
the matter we care about?  A neutron star's maximum mass, the luminosity
of an accretion disc, the thrust of a relativistic jet --- these all
depend on the specific form of $T^{\mu\nu}$ that enters the field
equations.  This section gives that form for two matter models: perfect
fluids and electromagnetic fields.  We also state the energy conditions
--- inequalities that ``reasonable'' matter should satisfy --- and note
where the codebase's wormhole spacetime violates them.

\paragraph{Physical content of the stress-energy tensor.}

The stress-energy tensor $T^{\mu\nu}$ encodes the density and flux of
energy and momentum.  In the rest frame of an observer with four-velocity
$\fourvel$, the components have direct physical interpretations: $T^{00}$
is the energy density (mass-energy per unit volume), $T^{0i}$ is the
momentum density (equivalently, the energy flux in the $i$-direction),
and $T^{ij}$ is the stress --- the flux of $i$-momentum across a surface
of constant $x^j$.  The diagonal spatial components $T^{ii}$ are
pressures; the off-diagonal components $T^{ij}$ ($i \neq j$) are shear
stresses.  The symmetry $T^{\mu\nu} = T^{\nu\mu}$ follows from the
requirement that angular momentum be conserved in the absence of
intrinsic spin --- the same symmetry that makes the Einstein tensor a
natural partner on the geometric side.
\physnote{The identification $T^{0i} = {}$momentum density $= {}$energy
flux is not a coincidence but a consequence of Lorentz invariance: the
requirement that $T^{\mu\nu}$ transform as a tensor under boosts forces
energy flux and momentum density to be components of the same object.}

\begin{figure}[htbp]
  \centering
  % stress-energy-components.tex — 4×4 grid showing T^{μν} component roles
% Inputable TikZ fragment for fig:stress-energy-components
% Requires: tikz with calc, positioning (loaded by ehbook.cls)
%           decorations.pathreplacing (loaded below)
%
\usetikzlibrary{decorations.pathreplacing}
\begin{tikzpicture}[
    font=\footnotesize,
  ]

  % ── Dimensions ──────────────────────────────────────────────────────
  \def\cw{1.35}   % cell width  (cm)
  \def\ch{1.10}   % cell height (cm)

  % Pre-compute key coordinates
  \pgfmathsetmacro{\gridW}{4*\cw}       % total grid width
  \pgfmathsetmacro{\gridH}{4*\ch}       % total grid height
  \pgfmathsetmacro{\braceX}{4*\cw+0.25} % brace x (grid right + gap)

  % ── Cell helper: \cell{col}{row}{colour}{opacity}{text} ────────────
  %    (col,row) = 0..3 with origin at top-left of the grid
  \newcommand{\cell}[5]{%
    \fill[#3, fill opacity=#4]
      ({#1*\cw},{-#2*\ch}) rectangle ({(#1+1)*\cw},{-(#2+1)*\ch});
    \node at ({(#1+0.5)*\cw},{-(#2+0.5)*\ch}) {#5};
  }

  % ── Colour fills and cell labels ───────────────────────────────────

  % Time-time: T^{00} — energy density (gold)
  \cell{0}{0}{EHAccretionGold}{0.15}{$T^{00}$}

  % Time-space row: T^{0i} — momentum density (crimson)
  \cell{1}{0}{EHHorizonCrimson}{0.10}{$T^{01}$}
  \cell{2}{0}{EHHorizonCrimson}{0.10}{$T^{02}$}
  \cell{3}{0}{EHHorizonCrimson}{0.10}{$T^{03}$}

  % Space-time column: T^{i0} — energy flux (crimson)
  \cell{0}{1}{EHHorizonCrimson}{0.10}{$T^{10}$}
  \cell{0}{2}{EHHorizonCrimson}{0.10}{$T^{20}$}
  \cell{0}{3}{EHHorizonCrimson}{0.10}{$T^{30}$}

  % Spatial diagonal: T^{ii} — pressure (blue)
  \cell{1}{1}{EHJetBlue}{0.15}{$T^{11}$}
  \cell{2}{2}{EHJetBlue}{0.15}{$T^{22}$}
  \cell{3}{3}{EHJetBlue}{0.15}{$T^{33}$}

  % Spatial off-diagonal: T^{ij}, i≠j — shear (gray)
  \cell{1}{2}{EHMarginGray}{0.10}{$T^{21}$}
  \cell{1}{3}{EHMarginGray}{0.10}{$T^{31}$}
  \cell{2}{1}{EHMarginGray}{0.10}{$T^{12}$}
  \cell{2}{3}{EHMarginGray}{0.10}{$T^{32}$}
  \cell{3}{1}{EHMarginGray}{0.10}{$T^{13}$}
  \cell{3}{2}{EHMarginGray}{0.10}{$T^{23}$}

  % ── Grid lines ─────────────────────────────────────────────────────

  % Thin grid lines
  \foreach \i in {0,...,4} {
    \draw[EHMarginGray, thin] ({\i*\cw},0) -- ({\i*\cw},{-\gridH});
    \draw[EHMarginGray, thin] (0,{-\i*\ch}) -- ({\gridW},{-\i*\ch});
  }

  % Heavier separator: 0-row/column vs 3×3 spatial block
  \draw[EHMarginGray, thick] ({\cw},0) -- ({\cw},{-\gridH});
  \draw[EHMarginGray, thick] (0,{-\ch}) -- ({\gridW},{-\ch});

  % Outer border
  \draw[EHMarginGray, semithick] (0,0) rectangle ({\gridW},{-\gridH});

  % ── Column headers (ν) ─────────────────────────────────────────────

  \node[above=4pt] at ({0.5*\cw},0) {$0\;(t)$};
  \node[above=4pt] at ({1.5*\cw},0) {$1\;(x)$};
  \node[above=4pt] at ({2.5*\cw},0) {$2\;(y)$};
  \node[above=4pt] at ({3.5*\cw},0) {$3\;(z)$};
  \node[above=16pt] at ({2*\cw},0) {$\nu$};

  % ── Row headers (μ) ────────────────────────────────────────────────

  \node[left=4pt] at (0,{-0.5*\ch}) {$0\;(t)$};
  \node[left=4pt] at (0,{-1.5*\ch}) {$1\;(x)$};
  \node[left=4pt] at (0,{-2.5*\ch}) {$2\;(y)$};
  \node[left=4pt] at (0,{-3.5*\ch}) {$3\;(z)$};
  \node[rotate=90, anchor=south] at (-1.2,{-2*\ch}) {$\mu$};

  % ── Right brace: spatial stress block ──────────────────────────────
  %   Path runs top-to-bottom; default brace opens left toward the grid.
  %   Label sits to the right of the brace tip.

  \draw[decorate, decoration={brace, amplitude=8pt},
        EHMarginGray, thick]
    (\braceX,{-\ch}) -- (\braceX,{-\gridH})
    node[midway, right=12pt, text=EHMarginGray, align=left]
      {\footnotesize spatial\\\footnotesize stress};

  % ── Legend ─────────────────────────────────────────────────────────

  \pgfmathsetmacro{\ly}{-\gridH - 0.55}  % legend row 1 y
  \pgfmathsetmacro{\lyb}{\ly - 0.4}      % legend row 2 y
  \def\sw{0.28}                            % swatch size (cm)
  \pgfmathsetmacro{\col}{2.7}             % second column x (cm)

  % Gold: energy density
  \fill[EHAccretionGold, fill opacity=0.35]
    (0,{\ly-\sw/2}) rectangle (\sw,{\ly+\sw/2});
  \draw[EHMarginGray, very thin]
    (0,{\ly-\sw/2}) rectangle (\sw,{\ly+\sw/2});
  \node[right=3pt, font=\scriptsize] at (\sw,\ly) {energy density};

  % Crimson: momentum / flux
  \fill[EHHorizonCrimson, fill opacity=0.25]
    (\col,{\ly-\sw/2}) rectangle ({\col+\sw},{\ly+\sw/2});
  \draw[EHMarginGray, very thin]
    (\col,{\ly-\sw/2}) rectangle ({\col+\sw},{\ly+\sw/2});
  \node[right=3pt, font=\scriptsize] at ({\col+\sw},\ly) {momentum / flux};

  % Blue: pressure
  \fill[EHJetBlue, fill opacity=0.35]
    (0,{\lyb-\sw/2}) rectangle (\sw,{\lyb+\sw/2});
  \draw[EHMarginGray, very thin]
    (0,{\lyb-\sw/2}) rectangle (\sw,{\lyb+\sw/2});
  \node[right=3pt, font=\scriptsize] at (\sw,\lyb) {pressure};

  % Gray: shear
  \fill[EHMarginGray, fill opacity=0.25]
    (\col,{\lyb-\sw/2}) rectangle ({\col+\sw},{\lyb+\sw/2});
  \draw[EHMarginGray, very thin]
    (\col,{\lyb-\sw/2}) rectangle ({\col+\sw},{\lyb+\sw/2});
  \node[right=3pt, font=\scriptsize] at ({\col+\sw},\lyb) {shear};

\end{tikzpicture}

  \caption[Physical content of the stress-energy tensor $T^{\mu\nu}$]{Physical content of the stress-energy tensor $T^{\mu\nu}$.
    The component $T^{00}$ (gold) is the energy density measured by an
    observer at rest.  The mixed time--space components $T^{0i}$ and
    $T^{i0}$ (crimson) give the momentum density and energy flux ---
    equal by the symmetry $T^{\mu\nu} = T^{\nu\mu}$.  Within the spatial
    $3\times 3$ block, the diagonal entries $T^{ii}$ (blue) are pressures
    and the off-diagonal entries $T^{ij}$ with $i \neq j$ (gray) are
    shear stresses.  For a perfect fluid, all shear components vanish and
    the three pressures are equal.}
  \label{fig:stress-energy-components}
\end{figure}

\paragraph{Perfect fluids.}

The simplest matter model with nontrivial dynamics is a perfect fluid:
matter with an energy density and an isotropic pressure, but no viscosity
or heat conduction.  In the rest frame of the fluid, the only nonzero
components are the energy density $\rho$ on the time--time diagonal and
the pressure $p$ on the spatial diagonal.  The covariant expression valid
in any frame is
%
\begin{equation}
  T^{\mu\nu} = (\rho + p)\,u^\mu\, u^\nu + p\, \inversemetric \,,
  \label{eq:perfect-fluid}
\end{equation}
%
where $u^\mu$ is the fluid four-velocity (normalised so that
$\metric\, u^\mu\, u^\nu = -1$), and $\rho$ and $p$ are measured in the
fluid's rest frame.  Here $\rho$ denotes the total energy density
(rest-mass plus internal energy); the rest-mass density $\restmassdensity$
used in the GRMHD formulation of Part~V is a distinct, smaller quantity.
The two terms in \cref{eq:perfect-fluid} have a clean separation: the
first carries the inertia (energy plus pressure, which acts as inertia in
relativity), directed along the fluid's worldline; the second is an
isotropic pressure that pushes equally in all spatial directions.
\physnote{In general relativity, pressure gravitates.  The combination
$\rho + p$ rather than $\rho$ alone appears because pressure contributes
to inertia --- a consequence of the equivalence of energy and mass.  This
has observable effects: it increases the central pressure required to
support a neutron star, ultimately setting a maximum mass (the
Tolman--Oppenheimer--Volkoff limit).}

The trace of \cref{eq:perfect-fluid} follows from contracting with
$\metric$ and using the normalisation $\metric\,u^\mu\,u^\nu = -1$
together with $\metric\,\inversemetric = 4$:
%
\begin{equation}
  T = \metric\,T^{\mu\nu} = -\rho + 3p \,.
  \label{eq:perfect-fluid-trace}
\end{equation}
%
The trace measures the balance between energy density and pressure: it
vanishes when $\rho = 3p$, the equation of state for a gas of massless
particles (a radiation-dominated fluid).  For dust --- pressureless matter
with $p = 0$ --- the stress-energy tensor reduces to
$T^{\mu\nu} = \rho\,u^\mu\,u^\nu$, and the trace gives $T = -\rho$.
Dust is the matter model behind the Newtonian limit discussed in
\cref{sec:gr-weak-field}: slow-moving, low-pressure matter whose
gravitational field is entirely determined by its mass density.

The conservation law $\nabla_\mu\,T^{\mu\nu} = 0$
(\cref{eq:stress-energy-conservation}) applied to a perfect fluid
yields the relativistic generalisation of the Euler equations --- a
continuity equation for energy and a momentum equation for the fluid ---
though the explicit decomposition belongs to the GRMHD treatment of
Part~V.
\xrefnote{The $3{+}1$ decomposition of perfect-fluid conservation and
the numerical flux formulation are developed in Part~V.}

\paragraph{The electromagnetic field.}

Magnetised accretion discs and relativistic jets are driven by
electromagnetic fields that carry energy and momentum --- and therefore
curve spacetime.  To couple these fields to gravity through the Einstein
equations, we need their stress-energy tensor.  The electromagnetic field
in curved spacetime is encoded in the antisymmetric Faraday tensor
$F_{\mu\nu} = -F_{\nu\mu}$, whose six independent components are the
electric and magnetic fields measured by a given observer.  The
stress-energy tensor built from $F_{\mu\nu}$ is
%
\begin{equation}
  T^{\mu\nu}_{\text{EM}}
    = F^{\mu}{}_{\alpha}\,F^{\nu\alpha}
    - \tfrac{1}{4}\,\inversemetric\,F_{\alpha\beta}\,F^{\alpha\beta} \,,
  \label{eq:em-stress-energy}
\end{equation}
%
where indices on $F$ are raised with the inverse metric.  Unlike the
perfect fluid, the electromagnetic stress-energy tensor is trace-free:
%
\begin{equation}
  \metric\,T^{\mu\nu}_{\text{EM}} = 0 \,.
  \label{eq:em-trace-free}
\end{equation}
%
Contracting $\metric$ with both terms in \cref{eq:em-stress-energy}: the
first yields $F_{\alpha\beta}\,F^{\alpha\beta}$, and the second, after
$\metric\,\inversemetric = 4$, yields
$-\tfrac{1}{4} \times 4 \times F_{\alpha\beta}\,F^{\alpha\beta}
= -F_{\alpha\beta}\,F^{\alpha\beta}$; the sum vanishes.  The vanishing
trace has an immediate consequence for the field equations: substituting
$T = 0$ into the trace-reversed form (\cref{eq:efe-trace-reversed})
gives $\ricci = 8\pi\,\stressenergy$, so the Ricci tensor is directly
proportional to the stress-energy tensor in electrovacuum.
\physnote{Maxwell's equations in curved spacetime take the same form as
in flat spacetime, with partial derivatives replaced by covariant
derivatives: $\nabla_\mu\,F^{\mu\nu} = 4\pi\,J^\nu$ (in Gaussian units)
and $\nabla_{[\alpha}\,F_{\mu\nu]} = 0$.  The stress-energy conservation
$\nabla_\mu\,T^{\mu\nu}_{\text{EM}} = 0$ then follows from Maxwell's
equations alone, independently of the Einstein equations.}

\paragraph{Energy conditions.}

Not every symmetric tensor deserves to sit on the right-hand side of the
field equations.  The energy conditions are physically motivated
inequalities that constrain what ``reasonable'' matter looks like.  The
two most important for this book are the weak and null energy conditions.
The \emph{weak energy condition} (WEC) requires that every observer
measure a non-negative energy density:
%
\begin{equation}
  T_{\mu\nu}\, v^\mu\, v^\nu \geq 0
  \qquad \text{for all timelike } v^\mu \,,
  \label{eq:wec}
\end{equation}
%
where $v^\mu$ is an arbitrary future-directed timelike vector.  For a
perfect fluid, one can verify that the minimum over all timelike
directions occurs in the rest frame (giving $\rho \geq 0$) and the
highly boosted limit adds $\rho + p \geq 0$.  The \emph{null energy
condition} (NEC) is the limiting case for null vectors:
%
\begin{equation}
  T_{\mu\nu}\, k^\mu\, k^\nu \geq 0
  \qquad \text{for all null } k^\mu \,.
  \label{eq:nec}
\end{equation}
%
For a perfect fluid, the null condition $\metric\,k^\mu\,k^\nu = 0$
kills the pressure term, leaving
$(\rho + p)(u_\mu\,k^\mu)^2 \geq 0$ --- so the NEC requires only
$\rho + p \geq 0$.  This is weaker than the WEC because null vectors
have no rest frame: the condition constrains only one combination of
$\rho$ and $p$, not both independently.  The NEC is the minimal
condition needed for Penrose's singularity theorem and for the area
theorem of classical black hole mechanics.

Two further conditions appear in the literature: the
\emph{strong energy condition} (SEC),
$(T_{\mu\nu} - \tfrac{1}{2}\,T\,g_{\mu\nu})\,v^\mu\,v^\nu \geq 0$ for
all timelike $v^\mu$, which ensures that gravity is attractive in the
sense of geodesic focusing; and the \emph{dominant energy condition}
(DEC), which adds to the WEC the requirement that the energy--momentum
current $-T^{\mu}{}_{\nu}\,v^\nu$ be causal (non-spacelike) for any
future-directed timelike $v^\mu$ --- physically, energy cannot flow faster
than light.  Ordinary matter satisfies all four: a perfect fluid obeys
them whenever $\rho \geq 0$ and $0 \leq p \leq \rho$, and the
electromagnetic field satisfies them for any field configuration.
\warnnote{The energy conditions are not laws of nature.  Quantum fields
can violate them locally (the Casimir effect produces regions of negative
energy density), and the cosmological constant violates the SEC.  They
are useful as sufficient conditions for theorems --- singularity
formation, area increase, topological censorship --- but exotic matter
that violates them is not forbidden by general relativity itself.}

\paragraph{Exotic matter and the Morris--Thorne wormhole.}

The energy conditions are abstract inequalities; a concrete violation
builds intuition for what they actually constrain.  The Morris--Thorne
wormhole implemented in the codebase (\cref{sec:gr-exact-solutions})
provides exactly this.  Its metric is static and spherically symmetric,
with a ``throat'' of minimum radius $b_0$ connecting two asymptotically
flat regions.  Sustaining this geometry requires a stress-energy tensor
that violates the NEC at the throat: any radial null vector $k^\mu$
satisfies $T_{\mu\nu}\,k^\mu\,k^\nu < 0$
there \cite{MorrisThorne1988}.  The violation is not a defect of the
solution but its defining feature --- without exotic matter, the throat
collapses to a singularity.  This is why the Morris--Thorne spacetime is
a pedagogical tool (and a visually striking one for the ray tracer) rather
than an astrophysically realistic solution.

\paragraph{From matter to geometry.}

The stress-energy tensors defined above complete the right-hand side of
the Einstein equations.  The spacetimes of this book fall into three
categories: vacuum solutions ($\stressenergy = 0$) for the Schwarzschild
and Kerr black holes, electrovacuum for magnetised discs, and
NEC-violating exotic matter for the Morris--Thorne wormhole.  With the
matter side of the field equations now concrete, the next section closes
the circle by verifying that the full equations reduce to Newtonian
gravity in the weak-field limit --- the most basic consistency check that
any relativistic theory of gravity must pass.
